% =====================================================================
% CHAPITRE 5 -- STANDALONE (compilable seul)
% =====================================================================
\documentclass[12pt,a4paper]{report}

% --- Encodage & langue ---
\usepackage[utf8]{inputenc}
\usepackage[T1]{fontenc}
\usepackage[french]{babel}

% --- Mise en page ---
\usepackage[top=2.5cm, bottom=2.5cm, left=2.5cm, right=2.5cm]{geometry}
\usepackage{float}
\usepackage{microtype}
\usepackage{parskip}

% --- Graphiques & TikZ ---
\usepackage{tikz}
\usetikzlibrary{shapes, arrows.meta, positioning, calc}

% --- Algorithmes ---
\usepackage[ruled,vlined,linesnumbered,french]{algorithm2e}
\SetKwInput{KwData}{\textcolor{primaryblue}{\textbf{Entrée}}}
\SetKwInput{KwResult}{\textcolor{scrumgreen}{\textbf{Sortie}}}
\newcommand\mycommfont[1]{\footnotesize\ttfamily\textcolor{gray}{#1}}
\SetCommentSty{mycommfont}

% --- Couleurs ---
\usepackage{xcolor}
\definecolor{primaryblue}{RGB}{0, 84, 166}
\definecolor{secondaryblue}{RGB}{0, 120, 215}
\definecolor{scrumgreen}{RGB}{0, 153, 76}
\definecolor{accentorange}{RGB}{230, 115, 0}

% --- Hyperliens ---
\usepackage{hyperref}
\hypersetup{colorlinks=true, linkcolor=primaryblue, urlcolor=primaryblue}

% --- Références fictives pour compilation autonome ---
% (remplacez par votre vrai fichier .bib si besoin)
\usepackage{filecontents}
\begin{filecontents*}{refs_ch5.bib}
@misc{sbert,
  author = {Reimers, Nils and Gurevych, Iryna},
  title  = {Sentence-BERT: Sentence Embeddings using Siamese BERT-Networks},
  year   = {2019},
  url    = {https://www.sbert.net}
}
@misc{react,
  author = {{Meta Platforms}},
  title  = {React -- A JavaScript library for building user interfaces},
  year   = {2023},
  url    = {https://react.dev}
}
\end{filecontents*}
\usepackage[backend=bibtex, style=numeric]{biblatex}
\addbibresource{refs_ch5.bib}

% =====================================================================
\begin{document}

% =====================================================================
% CHAPITRE 5
% =====================================================================

\chapter{Sprints 4 et 5 -- Intelligence Artificielle et Messagerie Avancée}

% ─────────────────────────────────────────────
\section*{Introduction}
% ─────────────────────────────────────────────

Ce chapitre couvre les deux derniers sprints du projet. Le Sprint~4 intègre un moteur de recommandation sémantique basé sur Sentence-BERT, une détection automatique des lacunes et un tuteur pédagogique Gemini, déployés sous forme d'un microservice Python FastAPI indépendant. Le Sprint~5 met en place un système de messagerie assurant la récupération sécurisée des comptes et l'envoi de notifications automatiques aux apprenants.


% ─────────────────────────────────────────────
\section{Sprint 4 -- Intelligence Artificielle et Recommandation Sémantique}
% ─────────────────────────────────────────────

Ce quatrième sprint, d'une durée de deux semaines, a été mené conjointement par ABDELHEDI Wassim pour le développement du moteur IA et FRIKHA Amin pour son intégration au backend. Son objectif principal consistait à concevoir un moteur de recommandation hybride exploitant le modèle Sentence-BERT, enrichi de mécanismes automatiques de détection des lacunes et d'un tuteur pédagogique interactif.

\subsection{Architecture du Service IA et Choix technologique : Sentence-BERT}

Le service de recommandation repose sur le modèle \texttt{paraphrase-multilingual-MiniLM-L12-v2} issu de la bibliothèque \textit{Sentence-Transformers}~\cite{sbert}. Il s'agit d'une architecture Transformer BERT fine-tunée pour la similarité sémantique, capable de produire des vecteurs de 384 dimensions couvrant plus de 50 langues. Son avantage majeur réside dans la recherche cross-lingue : une requête soumise en français peut retrouver des cours rédigés en anglais et vice-versa. Contrairement à une recherche classique par mots-clés, qui ne retourne que des correspondances exactes, Sentence-BERT évalue la proximité sémantique globale des énoncés, rendant la découverte de ressources beaucoup plus naturelle et pertinente pour l'apprenant.

\subsection{Sprint Backlog -- Sprint 4}

Les onze tâches planifiées pour ce sprint ont été réparties selon trois domaines de responsabilité. ABDELHEDI Wassim a pris en charge l'ensemble du développement du microservice IA : configuration de l'environnement FastAPI, module d'extraction textuelle multi-format, processeur NLP bilingue, moteur de recommandation hybride, détecteur de lacunes sémantique et service tuteur Gemini. FRIKHA Amin a assuré la couche d'intégration Spring Boot, comprenant les endpoints REST, le service de recherche WebClient asynchrone et la persistance de l'historique. BAHLOUL Donia a développé l'interface React avec la page de recherche multi-entités et l'affichage des recommandations dans le tableau de bord étudiant.

\begin{figure}[H]
\centering
\resizebox{0.88\textwidth}{!}{%
\begin{tikzpicture}[
  node distance=0.5cm and 1.4cm,
  member/.style={rectangle, rounded corners=4pt, draw=primaryblue, fill=primaryblue!12,
    font=\bfseries\small, minimum width=4cm, minimum height=0.72cm, align=center},
  task/.style={rectangle, rounded corners=3pt, draw=primaryblue!45, fill=white,
    font=\footnotesize, minimum width=4cm, minimum height=0.6cm, align=center,
    text width=3.8cm},
  arr/.style={->, thick, primaryblue!55}
]
\node[member] (w) {ABDELHEDI Wassim};
\node[task, below=of w] (t38) {Config. microservice FastAPI};
\node[task, below=of t38] (t39) {Extraction texte multi-format};
\node[task, below=of t39] (t40) {Processeur NLP bilingue};
\node[task, below=of t40] (t41) {Moteur recommandation Sentence-BERT};
\node[task, below=of t41] (t42) {Détecteur de lacunes};
\node[task, below=of t42] (t43) {Service tuteur Gemini};
\node[member, right=1.6cm of w] (a) {FRIKHA Amin};
\node[task, below=of a] (t44) {Endpoints REST + Pydantic};
\node[task, below=of t44] (t45) {SearchService WebClient};
\node[task, below=of t45] (t46) {Historique des recherches};
\node[member, right=1.6cm of a] (d) {BAHLOUL Donia};
\node[task, below=of d] (t47) {Page de recherche React};
\node[task, below=of t47] (t48) {Recommandations tableau de bord};
\draw[arr] (w)--(t38); \draw[arr] (t38)--(t39); \draw[arr] (t39)--(t40);
\draw[arr] (t40)--(t41); \draw[arr] (t41)--(t42); \draw[arr] (t42)--(t43);
\draw[arr] (a)--(t44); \draw[arr] (t44)--(t45); \draw[arr] (t45)--(t46);
\draw[arr] (d)--(t47); \draw[arr] (t47)--(t48);
\end{tikzpicture}}
\caption{Répartition des tâches -- Sprint Backlog Sprint~4}
\end{figure}

\section{Pipeline de Recommandation Hybride}

\subsection{Flux complet de recommandation}

Le pipeline de recommandation combine une analyse sémantique profonde à des ajustements contextuels pour affiner les résultats. Comme illustré dans le schéma ci-dessous, la requête de l'utilisateur traverse six étapes successives, de l'extraction des mots-clés jusqu'à la sélection finale des cours les plus pertinents.

\begin{figure}[H]
\centering
\begin{tikzpicture}[
  node distance=1.25cm,
  box/.style={
    rectangle, rounded corners=5pt,
    draw=primaryblue, fill=primaryblue!8,
    minimum width=9cm, minimum height=0.9cm,
    align=center, font=\small\sffamily
  },
  arrow/.style={->, thick, primaryblue!70, shorten >=2pt, shorten <=2pt}
]
\node[box] (q)   {\textbf{Requête utilisateur}};
\node[box, below=of q]   (kw)  {\textbf{Étape 1 — Extraction NLP}\\\footnotesize Analyse lexicale et identification des mots-clés};
\node[box, below=of kw]  (exp) {\textbf{Étape 2 — Expansion sémantique}\\\footnotesize Enrichissement par synonymes techniques bidirectionnels};
\node[box, below=of exp] (enc) {\textbf{Étape 3 — Encodage vectoriel}\\\footnotesize Génération d'embeddings de 384 dimensions via Sentence-BERT};
\node[box, below=of enc] (cos) {\textbf{Étape 4 — Similarité cosinus}\\\footnotesize Comparaison vectorielle avec le catalogue de cours indexé};
\node[box, below=of cos] (adj) {\textbf{Étape 5 — Scoring hybride}\\\footnotesize Pondérations : titre, catégorie, niveau et pénalité inscription};
\node[box, below=of adj] (top) {\textbf{Étape 6 — Agrégation et restitution}\\\footnotesize Top 10 cours + raison générée renvoyés au backend};
\draw[arrow] (q)--(kw); \draw[arrow] (kw)--(exp); \draw[arrow] (exp)--(enc);
\draw[arrow] (enc)--(cos); \draw[arrow] (cos)--(adj); \draw[arrow] (adj)--(top);
\end{tikzpicture}
\caption{Schéma explicatif du flux complet de recommandation hybride LearnAgent}
\end{figure}

\subsection{Mécanique de l'Algorithme de Scoring Hybride}

L'évaluation de la pertinence des cours repose sur une logique algorithmique en deux phases : une évaluation sémantique pure suivie de l'application de règles métiers (le scoring hybride).

\begin{algorithm}[H]
\caption{\textcolor{primaryblue}{\textbf{Algorithme de Recommandation Hybride (LearnAgent)}}}
\KwData{Requête $Q$, Liste des cours $C$, Cours inscrits $C_{inscrit}$, Top $N$}
\KwResult{Liste des $N$ cours recommandés}
$mots\_cles \leftarrow \text{ExtraireMotsCles}(Q)$\;
$Q_{etendu} \leftarrow \text{EnrichirSynonymes}(Q, mots\_cles)$\;
$E_Q \leftarrow \text{EncoderSentenceBERT}(Q_{etendu})$\;
$Resultats \leftarrow \emptyset$\;
\Pour{chaque cours $c \in C$}{
    $E_c \leftarrow \text{EncoderSentenceBERT}(c.texte)$\;
    $score \leftarrow \text{SimilariteCosinus}(E_Q, E_c)$\;
    \Si{$c.id \in C_{inscrit}$}{
        \textcolor{accentorange}{$score \leftarrow score \times 0.2$} \tcp*{Pénalité de découverte}
    }
    \Si{Titre de $c$ contient mots\_cles (exact)}{
        \textcolor{scrumgreen}{$score \leftarrow score + 0.15$}\;
    }
    \Si{Niveau de $c$ correspond à l'intention de $Q$}{
        \textcolor{scrumgreen}{$score \leftarrow score + 0.10$}\;
    } \Sinon {
        \textcolor{accentorange}{$score \leftarrow score - 0.05$}\;
    }
    \Si{$score \geq 0.10$}{
        Ajouter $(c, score)$ à $Resultats$\;
    }
}
Trier $Resultats$ par $score$ décroissant\;
\Retour{Les $N$ premiers éléments de $Resultats$}\;
\end{algorithm}

Le processus textuel s'organise ainsi :
\begin{enumerate}
    \item \textbf{Évaluation Sémantique} : L'algorithme compare le sens profond de la requête de l'apprenant avec le contenu textuel de chaque cours en calculant leur proximité vectorielle (Similarité Cosinus générée via Sentence-BERT).
    \item \textbf{Ajustements Heuristiques} : Le score de base de chaque cours est ensuite pondéré dynamiquement par plusieurs critères métier (Pénalité de découverte, Correspondance Directe, Adaptation au Niveau).
    \item \textbf{Filtrage et Restitution} : Un filtre de pertinence exclut systématiquement tous les cours dont le score final est inférieur à 0.1.
\end{enumerate}

\subsection{Expansion de Requête et NLP Bilingue}

Le module d'analyse linguistique enrichit automatiquement les requêtes courtes grâce à un dictionnaire de synonymes techniques bidirectionnels. Ainsi, un sigle ou un terme générique est étendu à ses variantes et acronymes associés pour maximiser les correspondances pertinentes. Ce traitement inclut également une détection heuristique de la langue de saisie, un filtrage des mots vides bilingue et une extraction de l'intention de l'utilisateur, permettant au moteur d'ajuster dynamiquement le score de pertinence selon le niveau de difficulté recherché.

\subsection{Pipeline d'Extraction de Texte et Mots-clés}

Lorsqu'un enseignant téléverse un support de cours, le microservice IA prend en charge l'extraction du contenu brut afin de l'indexer. Des parseurs spécifiques (\texttt{PyPDF2}, \texttt{python-docx}, \texttt{python-pptx}) extraient le texte des documents. Ensuite, le processeur NLP bilingue effectue un nettoyage en s'appuyant sur des dictionnaires de mots vides (stop-words) et des heuristiques linguistiques développées sur-mesure (sans dépendances lourdes). L'algorithme détecte la langue par un calcul de chevauchement ensembliste, puis sélectionne les termes techniques clés (max 20) en les dédupliquant et en conservant leur pertinence métier. Ces mots-clés sont retournés au backend Spring Boot pour enrichir les métadonnées du cours.

\begin{figure}[H]
\centering
\resizebox{0.9\textwidth}{!}{%
\begin{tikzpicture}[
  node distance=0.8cm and 1cm,
  doc/.style={rectangle, draw=gray, fill=gray!10, minimum width=2cm, minimum height=0.8cm, align=center, font=\footnotesize\ttfamily},
  proc/.style={rectangle, rounded corners=3pt, draw=primaryblue, fill=primaryblue!10, minimum width=3cm, minimum height=0.8cm, align=center, font=\footnotesize\sffamily},
  res/.style={rectangle, rounded corners=5pt, draw=scrumgreen, fill=scrumgreen!10, minimum width=3.5cm, minimum height=0.8cm, align=center, font=\small\bfseries},
  arrow/.style={->, thick, gray!80, >=stealth}
]
\node[doc] (pdf) {Document PDF};
\node[doc, below=of pdf] (docx) {Fichier DOCX};
\node[doc, below=of docx] (pptx) {Présentation PPTX};

\node[proc, right=1.5cm of docx] (extract) {Text Extractor\\(Service in-memory)};
\node[proc, right=1.5cm of extract] (nlp) {NLP Processor\\(Tokenisation, Stop-words)};
\node[res, right=1.5cm of nlp] (out) {Mots-clés \&\\Texte Nettoyé};

\draw[arrow] (pdf.east) -- (extract.west);
\draw[arrow] (docx.east) -- (extract.west);
\draw[arrow] (pptx.east) -- (extract.west);
\draw[arrow, primaryblue] (extract) -- (nlp) node[midway, above, font=\tiny, align=center] {Texte brut};
\draw[arrow, primaryblue] (nlp) -- (out) node[midway, above, font=\tiny, align=center] {JSON};
\end{tikzpicture}}
\caption{Processus d'extraction textuelle et d'analyse linguistique des cours}
\end{figure}

\subsection{Intégration Backend (SearchService)}

Le service de recherche du backend Spring Boot orchestre l'intégralité du flux transactionnel. Il agrège le catalogue des cours publiés et les données d'inscription de l'utilisateur, puis initie un appel asynchrone non-bloquant vers le microservice IA. Un mécanisme de secours automatique (fallback) bascule vers une recherche textuelle classique en base de données en cas d'inaccessibilité du moteur IA, garantissant ainsi la continuité du service. L'ensemble des requêtes est tracé dans un historique persistant, et les recommandations produites alimentent de manière proactive l'interface personnalisée de l'étudiant.

\section{Détection des Lacunes (Weak Topic Detector)}

\subsection{Algorithme et Règles de Détection}

La détection des lacunes intervient spécifiquement lorsqu'un apprenant échoue à une évaluation (score global inférieur à 60\%). 

\begin{algorithm}[H]
\caption{\textcolor{primaryblue}{\textbf{Algorithme de Détection des Lacunes Pédagogiques}}}
\KwData{Liste des questions échouées $Q_{err}$, Liste des chapitres $C$}
\KwResult{Rapport des lacunes avec niveaux de sévérité}
$Stats \leftarrow$ dictionnaire vide\;
\Pour{chaque question $q \in Q_{err}$}{
    \Si{$q$ n'a pas de chapitre assigné}{
        $E_q \leftarrow \text{EncoderSentenceBERT}(q.texte)$\;
        $chapitre\_proche \leftarrow \arg\max_{c \in C} \text{SimilariteCosinus}(E_q, E_c)$\;
        \Si{$\max(\text{Similarite}) > 0.3$}{
            \textcolor{scrumgreen}{$q.chapitre \leftarrow chapitre\_proche$}\;
        } \Sinon {
            \textcolor{accentorange}{$q.chapitre \leftarrow \text{"Concepts Généraux"}$}\;
        }
    }
    Incrémenter $Stats[q.chapitre].erreurs$\;
}
$Rapport \leftarrow \emptyset$\;
\Pour{chaque chapitre $c$ dans $Stats$}{
    $T_{erreur} \leftarrow \frac{Stats[c].erreurs}{\text{TotalQuestions}(c)} \times 100$\;
    \Si{$T_{erreur} > 70$}{
        Ajouter $(c, \text{\textcolor{red}{\textbf{Critique}}})$ à $Rapport$\;
    } \SinonSi{$T_{erreur} \geq 40$}{
        Ajouter $(c, \text{\textcolor{accentorange}{\textbf{Moyen}}})$ à $Rapport$\;
    } \Sinon {
        Ajouter $(c, \text{\textcolor{scrumgreen}{\textbf{Faible}}})$ à $Rapport$\;
    }
}
\Retour{$Rapport$}\;
\end{algorithm}

Le processus se résume ainsi :
\begin{enumerate}
    \item \textbf{Mapping Dynamique des Questions} : Si une question n'est pas explicitement rattachée à un chapitre, le système compare sémantiquement la question avec le contenu de chaque chapitre (Seuil $0.3$).
    \item \textbf{Agrégation des Erreurs} : L'algorithme calcule le pourcentage d'échec par concept.
    \item \textbf{Classification de Sévérité} : Le système catégorise l'urgence en Critique ($>70\%$), Moyen ($40\%-70\%$), ou Faible ($<40\%$).
\end{enumerate}

Le rapport généré compile les concepts \textbf{critiques} pour en alerter l'étudiant, en lui conseillant de réviser ces chapitres précis avant sa prochaine tentative.

\begin{figure}[H]
\centering
\begin{tikzpicture}[
  node distance=1.15cm and 1.8cm,
  step/.style={rectangle, rounded corners=4pt, draw=scrumgreen, fill=scrumgreen!10,
    minimum width=7.5cm, minimum height=0.85cm, align=center, font=\small\sffamily},
  level/.style={rectangle, rounded corners=3pt, draw=accentorange, fill=accentorange!10,
    minimum width=3cm, minimum height=0.75cm, align=center, font=\footnotesize\sffamily},
  arrow/.style={->, thick, scrumgreen!70},
  oarr/.style={->, thick, accentorange!70}
]
\node[step] (q1) {Échec au quiz — questions ratées identifiées};
\node[step, below=of q1] (q2) {Encodage Sentence-BERT des chapitres et des questions};
\node[step, below=of q2] (q3) {Comparaison cosinus : question $\rightarrow$ chapitre le plus proche};
\node[step, below=of q3] (q4) {Calcul du taux d'erreur par chapitre};
\node[step, below=of q4] (q5) {Classification et génération du rapport de lacunes};
\node[level, right=1.6cm of q3] (lc) {Critique $>$ 70\%};
\node[level, above=0.5cm of lc] (ll) {Faible $<$ 40\%};
\node[level, below=0.5cm of lc] (lm) {Moyen 40--70\%};
\draw[arrow] (q1)--(q2); \draw[arrow] (q2)--(q3);
\draw[arrow] (q3)--(q4); \draw[arrow] (q4)--(q5);
\draw[oarr] (q5.east) -- ++(0.7,0) |- (ll);
\draw[oarr] (q5.east) -- ++(0.7,0) |- (lc);
\draw[oarr] (q5.east) -- ++(0.7,0) |- (lm);
\end{tikzpicture}
\caption{Schéma du processus de détection automatique des lacunes pédagogiques}
\end{figure}

\section{Tuteur Pédagogique IA (Gemini)}

Afin de proposer un retour formatif personnalisé, un tuteur intelligent basé sur l'API Google Gemini complète le dispositif de détection des lacunes. Lorsqu'une erreur est identifiée, le service transmet la question, la réponse de l'étudiant et la correction attendue au modèle de langage. Ce dernier génère une explication concise et constructive, guidant l'apprenant vers la compréhension du concept sans se limiter à lui livrer la réponse exacte. 

Pour optimiser les performances et la résilience, le service emploie deux stratégies techniques majeures :
\begin{itemize}
    \item \textbf{Traitement par lots (Batch Processing)} : L'ensemble des questions échouées est concaténé dans un seul \textit{prompt} structuré demandant explicitement un retour sous format \texttt{JSON}. Cela réduit l'empreinte réseau d'un facteur $O(n)$ à $O(1)$.
    \item \textbf{Chaîne de secours (Fallback Routing)} : En cas d'erreur de quota (HTTP 429) ou d'indisponibilité, le service bascule automatiquement à travers un tableau de modèles successifs (ex. \texttt{gemini-2.0-flash}, \texttt{gemini-pro-latest}) garantissant une disponibilité continue à l'utilisateur.
\end{itemize}

\begin{figure}[H]
\centering
\resizebox{0.85\textwidth}{!}{%
\begin{tikzpicture}[
  node distance=1.2cm and 1.5cm,
  action/.style={rectangle, rounded corners=4pt, draw=accentorange, fill=accentorange!12,
    minimum width=3.5cm, minimum height=0.8cm, align=center, font=\small\sffamily},
  model/.style={ellipse, draw=primaryblue, fill=primaryblue!10,
    minimum width=2.5cm, minimum height=1cm, align=center, font=\small\bfseries},
  arrow/.style={->, thick, primaryblue!70, >=stealth}
]
\node[action] (fail) {Questions échouées\\(Texte, Rép. étudiant, Vraie rép.)};
\node[action, right=of fail] (batch) {Agrégation en Batch\\(Optimisation temps de réponse)};
\node[model, below=of batch] (gemini) {Modèle\\Google Gemini};
\node[action, left=of gemini] (feedback) {Génération d'explications\\pédagogiques ciblées};
\node[action, below=of feedback] (ui) {Affichage UI :\\« Tuteur IA »};

\draw[arrow] (fail) -- (batch);
\draw[arrow] (batch) -- (gemini) node[midway, right, font=\footnotesize] {Prompt structuré};
\draw[arrow] (gemini) -- (feedback) node[midway, above, font=\footnotesize] {JSON Responses};
\draw[arrow] (feedback) -- (ui);
\end{tikzpicture}}
\caption{Flux d'intégration du Tuteur IA Gemini pour la correction de quiz}
\end{figure}

\section{Endpoints REST du Service IA}

Le microservice FastAPI expose une interface RESTful structurée autour des fonctionnalités clés de la couche IA. Le schéma suivant présente les principaux points d'accès et leur rôle respectif au sein de la plateforme.

\begin{figure}[H]
\centering
\resizebox{0.92\textwidth}{!}{%
\begin{tikzpicture}[
  node distance=0.65cm and 0.5cm,
  ep/.style={rectangle, rounded corners=4pt, draw=primaryblue, fill=primaryblue!10,
    minimum width=5.5cm, minimum height=0.72cm, align=center, font=\footnotesize\ttfamily},
  desc/.style={rectangle, rounded corners=3pt, draw=gray!40, fill=gray!8,
    minimum width=6.2cm, minimum height=0.72cm, align=left, font=\footnotesize\sffamily,
    text width=6cm, inner sep=4pt},
  arr/.style={->, thick, primaryblue!60}
]
\node[ep] (e1) {POST /api/recommend};
\node[desc, right=1cm of e1] (d1) {Recommandation hybride Sentence-BERT};
\node[ep, below=of e1] (e2) {POST /api/index-course};
\node[desc, right=1cm of e2] (d2) {Indexation du cours + extraction NLP};
\node[ep, below=of e2] (e3) {POST /api/extract-text};
\node[desc, right=1cm of e3] (d3) {Extraction de texte : PDF, DOCX, PPTX, TXT};
\node[ep, below=of e3] (e4) {POST /api/detect-weak-topics};
\node[desc, right=1cm of e4] (d4) {Détection sémantique des lacunes après quiz};
\node[ep, below=of e4] (e5) {POST /api/batch-tutor-feedback};
\node[desc, right=1cm of e5] (d5) {Feedback pédagogique groupé via Gemini};
\node[ep, below=of e5] (e6) {GET /health};
\node[desc, right=1cm of e6] (d6) {Vérification de disponibilité du microservice};
\foreach \s/\t in {e1/d1,e2/d2,e3/d3,e4/d4,e5/d5,e6/d6} \draw[arr] (\s)--(\t);
\end{tikzpicture}}
\caption{Endpoints REST du microservice IA FastAPI}
\end{figure}

\section{Architecture du Moteur IA (Python FastAPI)}

\subsection{Structure du service}

Le moteur de recommandation est un microservice Python indépendant (\texttt{ai-recommendation-engine/}), dont l'architecture est présentée à la figure~\ref{fig:ai_archi}. Il est composé des modules suivants :

\begin{itemize}
  \item \textbf{main.py} : Point d'entrée FastAPI avec configuration CORS et enregistrement des routes.
  \item \textbf{models/recommender.py} : Classe \texttt{CourseRecommender} -- chargement Sentence-BERT, pipeline d'embedding et de recommandation.
  \item \textbf{models/nlp\_processor.py} : \texttt{NLPProcessor} -- extraction de mots-clés, détection de langue (FR/EN).
  \item \textbf{services/text\_extractor.py} : Extraction de texte PDF (PyPDF2), DOCX (python-docx), PPTX (python-pptx).
  \item \textbf{services/weak\_topic\_detector.py} : Détection des lacunes par similarité sémantique entre questions ratées et chapitres.
  \item \textbf{routes/recommendation\_routes.py} : 6 endpoints REST exposés au backend Spring Boot.
  \item \textbf{schemas/recommendation\_schema.py} : Modèles Pydantic pour la validation des requêtes/réponses.
\end{itemize}

\begin{figure}[H]
\centering
\begin{tikzpicture}[
  mod/.style={rectangle, draw=scrumgreen, rounded corners=4pt, fill=scrumgreen!10, minimum width=3.8cm, minimum height=0.85cm, align=center, font=\small},
  arr/.style={->, thick, >=Stealth, color=scrumgreen}
]
\node[mod, fill=scrumgreen!25] (main) at (0,0)      {main.py\\(FastAPI)};
\node[mod] (routes)   at (-4.5,-1.8)  {recommendation\_routes.py\\(5 endpoints REST)};
\node[mod] (reco)     at (0,-1.8)     {recommender.py\\(Sentence-BERT)};
\node[mod] (nlp)      at (4.5,-1.8)   {nlp\_processor.py\\(Extraction mots-clés)};
\node[mod] (extract)  at (-2.5,-3.5)  {text\_extractor.py\\(PDF/DOCX/PPTX)};
\node[mod] (weak)     at (2.5,-3.5)   {weak\_topic\_detector.py\\(Lacunes IA)};

\draw[arr] (main) -- (routes);
\draw[arr] (main) -- (reco);
\draw[arr] (main) -- (nlp);
\draw[arr] (reco) -- (extract);
\draw[arr] (reco) -- (weak);
\end{tikzpicture}
\caption{Architecture du microservice Python FastAPI (moteur IA)}
\label{fig:ai_archi}
\end{figure}

\section{Sprint Review -- Sprint 4}

La clôture du Sprint~4 valide l'intégration réussie d'une couche d'intelligence artificielle autonome au sein de la plateforme. Le moteur de recommandation hybride a démontré sa capacité à relier sémantiquement des requêtes multilingues à un catalogue de formations diversifié, enrichi par une expansion lexicale bidirectionnelle. La détection automatique des lacunes après échec à un quiz et le tuteur pédagogique Gemini enrichissent considérablement l'expérience d'apprentissage personnalisé. Sur le plan architectural, l'isolation du microservice IA garantit des performances stables et une résilience accrue grâce au mécanisme de fallback, validant ainsi la pertinence d'une approche microservices et d'une communication inter-systèmes asynchrone.

\newpage

% ─────────────────────────────────────────────
% ─────────────────────────────────────────────
\section{Sprint 5 -- Système de Messagerie et Notifications}
% ─────────────────────────────────────────────

Ce cinquième et dernier sprint a été mené par FRIKHA Amin pour le développement backend et BAHLOUL Donia pour la conception des interfaces. Son objectif visait à mettre en place un écosystème de communication complet, intégrant d'une part l'envoi de notifications par courriel (récupération de comptes, alertes de mise à jour), et d'autre part une messagerie interne instantanée pour faciliter les échanges directs entre étudiants et enseignants.

\subsection{Architecture du module de messagerie et de notifications}

Le dispositif se divise en deux volets majeurs :
\begin{itemize}
    \item \textbf{Notifications asynchrones (Email)} : Un service d'envoi d'emails adossé au protocole SMTP génère des courriels dynamiques via des templates HTML (ex. \texttt{lesson-updated.html}). Il orchestre la création de jetons temporaires pour la réinitialisation de mot de passe (\texttt{PasswordResetToken}) et utilise des tâches planifiées (\texttt{EmailScheduler}) pour relancer les étudiants inactifs.
    \item \textbf{Messagerie interne (In-App Chat)} : Un contrôleur dédié (\texttt{MessageController}) gère les conversations entre utilisateurs. Il permet la recherche de contacts, le suivi de l'historique des discussions et la comptabilisation des messages non lus.
\end{itemize}

\begin{figure}[H]
\centering
\resizebox{0.9\textwidth}{!}{%
\begin{tikzpicture}[
  node distance=1cm and 1.5cm,
  comp/.style={rectangle, rounded corners=5pt, draw=secondaryblue, fill=secondaryblue!10,
    minimum width=4cm, minimum height=0.8cm, align=center, font=\small\sffamily},
  event/.style={rectangle, rounded corners=4pt, draw=accentorange, fill=accentorange!10,
    minimum width=4.5cm, minimum height=0.7cm, align=center, font=\footnotesize\sffamily},
  arrow/.style={->, thick, secondaryblue!60}
]
% Colonne Email
\node[comp] (smtp) {Service SMTP (Email)};
\node[comp, below=of smtp] (sched) {EmailScheduler};
\node[event, right=1cm of smtp] (ev1) {Jetons \& Mots de passe};
\node[event, right=1cm of sched] (ev2) {Alertes (inactivité, maj)};

% Colonne In-App
\node[comp, below=1.5cm of sched] (chat) {MessageController (In-App)};
\node[event, right=1cm of chat] (ev3) {Conversations \& Historique};
\node[event, below=0.3cm of ev3] (ev4) {Recherche de contacts};

\draw[arrow] (smtp)--(ev1); 
\draw[arrow] (sched)--(ev2); 
\draw[arrow] (chat)--(ev3);
\draw[arrow] (chat)--(ev4);
\end{tikzpicture}}
\caption{Architecture globale des services de communication (Sprint~5)}
\end{figure}

\subsection{Intégration Frontend}

Côté client, l'expérience utilisateur a été enrichie de nouvelles interfaces ergonomiques dédiées à la gestion de l'identité numérique. Les flux de réinitialisation de mot de passe s'intègrent naturellement au design system existant, permettant à l'utilisateur de récupérer ses accès de manière autonome, intuitive et totalement sécurisée, sans compromettre l'architecture d'authentification JWT existante.

% ─────────────────────────────────────────────
\section*{Conclusion}
\addcontentsline{toc}{section}{Conclusion}
% ─────────────────────────────────────────────

Ce chapitre a détaillé la réalisation des deux ultimes sprints du projet, scellant l'aboutissement technique de la plateforme LearnAgent. L'intégration du moteur de recommandation sémantique Sentence-BERT et du tuteur pédagogique Gemini dote la plateforme d'une réelle capacité d'adaptation aux besoins de chaque apprenant. Associée au système de notification intelligent du Sprint~5, la solution offre désormais une couverture fonctionnelle globale : de la diffusion du savoir à l'évaluation personnalisée, en passant par le suivi individualisé et l'engagement continu des étudiants.

\newpage

% =====================================================================
% CONCLUSION GÉNÉRALE
% =====================================================================
\chapter*{Conclusion Générale et Perspectives}
\addcontentsline{toc}{chapter}{Conclusion Générale et Perspectives}
\markboth{Conclusion Générale et Perspectives}{}

Au terme de ce Projet de Fin d'Année, nous avons conçu et développé \textbf{LearnAgent}, une plateforme e-learning intelligente et complète répondant à l'ensemble des besoins identifiés lors de l'analyse de l'existant. Structuré en cinq sprints Scrum successifs, le projet a permis de livrer de façon incrémentielle un système d'authentification JWT sécurisé avec gestion des rôles, un moteur de cours multi-chapitres avec progression stricte, des modules interactifs de quiz et d'exercices pratiques, puis un moteur de recommandation sémantique Sentence-BERT couplé à un tuteur pédagogique Gemini, et enfin un système robuste de messagerie et de notifications automatiques.

Sur le plan des compétences, ce projet nous a permis de consolider notre maîtrise des architectures microservices et de la conception d'API RESTful, en combinant Spring Boot 3.2 pour le backend transactionnel, React 19~\cite{react} pour le frontend et Python FastAPI pour la couche d'intelligence artificielle. Nous avons approfondi la gestion de la sécurité applicative (JWT, BCrypt, RBAC), le traitement du langage naturel multilingue avec Sentence-BERT et la similarité cosinus, ainsi que l'intégration des grands modèles de langage via l'API Gemini. Au-delà des aspects techniques, ce projet nous a également permis de renforcer notre capacité à travailler en équipe, à gérer les priorités dans un cadre agile et à documenter rigoureusement nos travaux.

En perspective, plusieurs axes d'amélioration se dégagent naturellement de ce premier prototype fonctionnel. L'intégration de modèles de langage de grande taille (GPT-4, Gemini) permettrait de générer des contenus pédagogiques personnalisés et d'enrichir le tuteur virtuel. Un moteur d'apprentissage adaptatif dynamique, ajustant le parcours de chaque étudiant en temps réel selon ses performances, constituerait une avancée significative. Par ailleurs, le déploiement d'une application mobile React Native, l'ajout de fonctionnalités sociales (forums, sessions collaboratives) et la mise en place de mécanismes de gamification (badges, classements, points de progression) contribueraient à renforcer l'engagement des apprenants. Enfin, la containerisation Docker et l'orchestration Kubernetes sur un hébergement cloud garantiraient la scalabilité et la résilience de la plateforme à grande échelle.

\vspace{1cm}
\begin{flushright}
\textit{BAHLOUL Donia, FRIKHA Amin, ABDELHEDI Wassim}\\
\textit{Sfax, 2026}
\end{flushright}

\newpage

% =====================================================================
% BIBLIOGRAPHIE
% =====================================================================
\printbibliography[heading=bibintoc, title={Bibliographie}]

\end{document}